% =====================================================================
% CHAPTER 2: Literature Review  (~24 pages, the masterclass chapter)
% WRITING SESSION INSTRUCTIONS (keep until the chapter is done):
%   - Read thesis/WRITING_PLAN.md section 4 FIRST. Every subsection there
%     names the exact summary files to read BEFORE writing it.
%   - Invoke ruet-thesis-expert + research-linguistics-expert. Register:
%     flagship for background, "[SYSTEM] [ref] [verb]s [mechanism]" for
%     minor works, fuller treatment for the ~10 anchor architectures.
%   - The stereo-vision-expert agent arbitrates ambiguous claims.
%     Numbers ONLY from papers/verified_performance.md (per-cell paper/
%     table/page breadcrumbs) or method_data.py.
%   - Figures and tables below are DONE (embedded, cited); do not move
%     them; write the prose around them.
%   - Architecture panels in Fig \ref{fig:paradigms} are REPRODUCED from
%     the original papers with citation (user-approved practice for the
%     thesis); all other diagrams are our own redrawn .drawio figures.
%   - PO demonstrated: PO2 throughout.
% =====================================================================
\chapter{Literature Review}

\section{Stereo Matching Fundamentals}
% Sources: tier1/surveys/Scharstein_Taxonomy_IJCV2002.md,
% tier1/surveys/Hirschmuller_SGM_TPAMI2007.md, papers/CONCEPTS.md,
% review_paper/sections/02_background.tex.
\todo{Write per WRITING\_PLAN section 4.1 row 2.1: rectification,
epipolar constraint, disparity; the four-step taxonomy
\cite{scharstein2002taxonomy}; classical costs (SAD, NCC, census);
SGM \cite{hirschmuller2008sgm} as the classical apex.}

\begin{figure}[!htbp]
\centering
\includegraphics[width=\textwidth]{fig_2_1_classical_preview}
\caption{\textbf{The classical stereo pipeline.} The four canonical steps of the Scharstein and Szeliski taxonomy \cite{scharstein2002taxonomy}: matching cost computation, cost aggregation, disparity optimization, and disparity refinement.}
\label{fig:classical_pipeline}
\end{figure}

\todo{Eq 2.1 SGM energy $E(D)$ with $P_1$/$P_2$ terms
\cite{hirschmuller2008sgm}; label eq:sgm\_energy.}

\section{The Deep Learning Transition}
% Sources: tier2/end_to_end/{MC-CNN,DispNetC,GCNet}.md,
% tier1/surveys/Poggi_Synergies_TPAMI2021.md.
\todo{Write per WRITING\_PLAN section 4.1 row 2.2: MC-CNN
\cite{zbontar2016mccnn} (learned cost, classical rest); DispNetC
\cite{mayer2016sceneflow} (first end-to-end, correlation layer); GC-Net
\cite{kendall2017gcnet} (4D cost volume + 3D conv + soft-argmin).}

\todo{Eq 2.2 correlation layer \cite{mayer2016sceneflow}; label
eq:correlation.}
\todo{Eq 2.3 soft-argmin regression \cite{kendall2017gcnet}; label
eq:softargmin.}

\section{The 3D Cost-Volume Era}
% Sources: tier2/end_to_end/{PSMNet,GWCNet,GANet,AANet,ACVNet,CFNet}.md.
\todo{Write per WRITING\_PLAN section 4.1 row 2.3: PSMNet
\cite{chang2018psmnet}; GwcNet group-wise correlation
\cite{guo2019gwcnet} (our tile initialization adopts this op, say so);
GA-Net \cite{zhang2019ganet}; one clause each for AANet
\cite{xu2020aanet}, ACVNet \cite{xu2022acvnet}, CFNet \cite{shen2021cfnet}.}

\todo{Eq 2.4 group-wise correlation \cite{guo2019gwcnet}; label
eq:gwc\_lit; forward-reference our \eqref{eq:gwc}.}

\section{Iterative Refinement}
% Sources: tier1/iterative/* + reference_impls/raft_stereo, igev_stereo.
\todo{Write per WRITING\_PLAN section 4.1 row 2.4: RAFT-Stereo
\cite{lipson2021raftstereo} (all-pairs correlation + ConvGRU + convex
upsampling: three mechanisms our model inherits in modified form);
IGEV \cite{xu2023igev} (Geometry Encoding Volume: direct ancestor of our
GEV block); CREStereo \cite{li2022crestereo}; Selective-Stereo
\cite{wang2024selective}.}

\todo{Eq 2.5 ConvGRU update \cite{lipson2021raftstereo}; label
eq:convgru\_lit.}
\todo{Eq 2.6 Geometry Encoding Volume \cite{xu2023igev}; label
eq:gev\_lit.}

\begin{figure}[!htbp]
\centering
\begin{subfigure}{0.49\textwidth}
\centering
\includegraphics[width=\textwidth,height=3.4cm,keepaspectratio]{fig_2_4a_psmnet}
\caption{PSMNet: 3D cost-volume paradigm \cite{chang2018psmnet}}
\end{subfigure}\hfill
\begin{subfigure}{0.49\textwidth}
\centering
\includegraphics[width=\textwidth,height=3.4cm,keepaspectratio]{fig_2_4b_raftstereo}
\caption{RAFT-Stereo: iterative paradigm \cite{lipson2021raftstereo}}
\end{subfigure}\\[6pt]
\begin{subfigure}{0.49\textwidth}
\centering
\includegraphics[width=\textwidth,height=3.4cm,keepaspectratio]{fig_2_4c_hitnet}
\caption{HITNet: tile-hypothesis paradigm \cite{tankovich2021hitnet}}
\end{subfigure}\hfill
\begin{subfigure}{0.49\textwidth}
\centering
\includegraphics[width=\textwidth,height=3.4cm,keepaspectratio]{fig_2_4d_liteanystereo}
\caption{LiteAnyStereo: efficient KD-trained paradigm \cite{jing2025liteanystereo}}
\end{subfigure}
\caption{\textbf{Four stereo-matching paradigms.} Architecture overviews reproduced from the original publications for comparison; each caption cites the source. The proposed StereoLite combines the tile-hypothesis state of (c) with the recurrent refinement of (b) and a geometry encoding volume in the spirit of IGEV.}
\label{fig:paradigms}
\end{figure}

\section{Foundation-Model Era}
% Sources: tier1/foundation_model/* + review_paper section 03.
\todo{Write per WRITING\_PLAN section 4.1 row 2.5: FoundationStereo
\cite{wen2025foundationstereo}, DEFOM-Stereo \cite{jiang2025defom},
MonSter \cite{cheng2025monster}, Stereo-Anywhere
\cite{bartolomei2025stereoanywhere}; the common recipe (monocular prior
+ massive data); why these motivate rather than solve the edge problem.}

\begin{figure}[!htbp]
\centering
\includegraphics[width=\textwidth]{fig_2_2_timeline_preview}
\caption{\textbf{Evolution of deep stereo matching, 2016 to 2026.} Representative methods arranged by publication year within six paradigm lanes; the outlined marker places this thesis relative to the modern efficient family surveyed in Section 2.6.}
\label{fig:timeline}
\end{figure}

\section{Efficient and Edge-Deployable Stereo}
% THE load-bearing section (~7 pages). Skeleton: review_paper section 04
% (rewrite for thesis register, do not copy).
\todo{Family-by-family prose per WRITING\_PLAN section 4.1 row 2.6,
following the tree of Fig \ref{fig:taxonomy}.}

\begin{figure}[!htbp]
\centering
\includegraphics[width=\textwidth]{fig_2_5_taxonomy_preview}
\caption{\textbf{Taxonomy of compression techniques for deep stereo matching.} Seven technique families organized by the pipeline stage they compress, each with representative methods; this tree organizes the survey of Section 2.6.}
\label{fig:taxonomy}
\end{figure}

\subsection{Backbone Substitution}
\todo{MobileStereoNet \cite{shamsafar2022mobilestereonet};
lightweight encoders.}
\subsection{Cost-Volume Compression}
\todo{cascade volumes \cite{gu2020cascade}, bilateral grid
\cite{xu2021bgnet}, pruning \cite{duggal2019deeppruner}, tile
hypotheses \cite{tankovich2021hitnet} (our lineage, say so).}
\subsection{Iterative-Loop Compression}
\todo{LightStereo \cite{guo2025lightstereo}, DTP \cite{pan2024dtp}.}
\subsection{Knowledge Distillation}
\todo{LiteAnyStereo three-stage KD \cite{jing2025liteanystereo}
(anchor; decisive for our future work), Distill-then-Prune.}
\todo{Eq 2.7 teacher-student disparity distillation loss
\cite{jing2025liteanystereo}; label eq:kd\_loss.}

\begin{figure}[!htbp]
\centering
\includegraphics[width=\textwidth]{fig_2_7_kd_stages}
\caption{\textbf{The three-stage knowledge-distillation pipeline of LiteAnyStereo.} Reproduced from \cite{jing2025liteanystereo}: supervised synthetic training, self-distillation with perturbation, and fine-tuning on unlabeled real data with pseudo labels from a frozen accurate model.}
\label{fig:kd_stages}
\end{figure}

\subsection{Architectural Compression and Adaptive Compute}
\todo{CoEx guided excitation \cite{bangunharcana2021coex}, CGI-Stereo
\cite{xu2023cgistereo}; AnyNet anytime inference \cite{wang2019anynet},
MADNet online adaptation \cite{tonioni2019madnet}; NAS: LEAStereo
\cite{cheng2020leastereo} in one paragraph.}

\section{Cross-Domain Generalization}
% Sources: tier3/domain_shift/*, review_paper section 05.
\todo{Write per WRITING\_PLAN section 4.1 row 2.7: why synthetic-trained
models fail on real domains; architecture-level remedies (DSMNet
\cite{zhang2020dsmnet}, HVT \cite{chang2023hvt}) vs training-recipe
remedies (StereoAnything \cite{guo2024stereoanything}, LiteAnyStereo
\cite{jing2025liteanystereo}); the evidence that the training pipeline,
not iteration, carries generalization.}

\section{Benchmarks and Datasets}
% Sources: tier2/datasets/*, review_paper section 02.3.
\todo{Prose around Table \ref{tab:datasets} per WRITING\_PLAN section 4.1
row 2.8, including the Middlebury eval3 protocol (15 training pairs with
public GT + 15 test pairs with hidden GT; the zero-shot convention
evaluates on the GT-public scenes).}

\begin{table}[!htbp]
\centering
\caption{\textbf{Stereo benchmarks and training datasets.} Pair counts are the standard training/test splits; SceneFlow counts follow the canonical finalpass split used in this thesis.}
\label{tab:datasets}
\small
\begin{tabular}{lllllc}
\toprule
Dataset & Domain & Train/Test & Resolution & Ground truth & Ref.\\
\midrule
SceneFlow & synthetic & 35{,}454 / 4{,}370 & $960{\times}540$ & dense, rendered & \cite{mayer2016sceneflow}\\
KITTI 2012 & driving & 194 / 195 & $1242{\times}375$ & sparse LiDAR & \cite{geiger2012kitti}\\
KITTI 2015 & driving & 200 / 200 & $1242{\times}375$ & sparse LiDAR + CAD & \cite{menze2015kitti}\\
Middlebury 2014 & indoor & 15 / 15 & up to ${\sim}3000{\times}2000$ & structured light & \cite{scharstein2014middlebury}\\
ETH3D & mixed & 27 / 20 & ${\sim}940{\times}490$ & laser scan & \cite{schops2017eth3d}\\
\bottomrule
\end{tabular}
\end{table}

\section{Summary and Research Gap}
\todo{Synthesis prose around Table \ref{tab:method_comparison} per
WRITING\_PLAN section 4.1 row 2.9; the gap statement uses the pattern
"the surveyed literature does not directly address [sub-3 M tile+GEV
fusion, real-time on Orin-class hardware, robustness to imperfect
rectification]". A review summarises; no roadmaps.}

\begin{table}[!htbp]
\centering
\caption{\textbf{Representative methods across eras.} Parameters, SceneFlow end-point error, KITTI 2015 D1-all, and reported latency; every value verified against the cited paper's tables (latency hardware differs across papers and is not normalized). The StereoLite row reports our full FlyingThings3D test result at the native axis (Chapter 4) and RTX 3050 fp16 latency.}
\label{tab:method_comparison}
\small
\begin{tabular}{lccccc}
\toprule
Method & Year & Params (M) & SF EPE (px) & KITTI D1 (\%) & Latency (ms)\\
\midrule
PSMNet \cite{chang2018psmnet} & 2018 & 5.2 & 1.09 & 2.32 & 410\\
HITNet (L) \cite{tankovich2021hitnet} & 2021 & 0.97 & 0.43 & 1.98 & 54\\
BGNet \cite{xu2021bgnet} & 2021 & ${\sim}2.9$ & 1.17 & 2.51 & 25\\
CoEx \cite{bangunharcana2021coex} & 2021 & 2.72 & 0.69 & 2.13 & 27\\
RAFT-Stereo \cite{lipson2021raftstereo} & 2021 & 11.23 & 0.61 & 1.82 & 380\\
IGEV-Stereo \cite{xu2023igev} & 2023 & 12.60 & 0.47 & 1.59 & 180\\
LightStereo-S \cite{guo2025lightstereo} & 2025 & 3.44 & 0.73 & 2.30 & 17\\
DEFOM-Stereo \cite{jiang2025defom} & 2025 & 47.3 & 0.42 & 1.55 & 316\\
FoundationStereo \cite{wen2025foundationstereo} & 2025 & ${\sim}340$ & 0.34 & 1.46 & ${\sim}470$\\
\midrule
StereoLite (ours) & 2026 & 2.96 & 0.78 & — & 49.8\\
\bottomrule
\end{tabular}
\end{table}

\begin{figure}[!htbp]
\centering
\begin{subfigure}{\textwidth}
\centering
\includegraphics[width=0.94\textwidth]{fig_2_6a_param_pareto}
\caption{accuracy against trainable parameters}
\end{subfigure}\\[6pt]
\begin{subfigure}{\textwidth}
\centering
\includegraphics[width=0.94\textwidth]{fig_2_6b_latency_pareto}
\caption{KITTI 2015 D1-all against reported latency}
\end{subfigure}
\caption{\textbf{Accuracy against parameters and latency across the surveyed literature.} Marker shape encodes the method family; every value carries a per-cell source citation in the underlying data table.}
\label{fig:pareto}
\end{figure}
